\section{从 Online Robot Learning 到 Offline Recipe}

这一章把三种原料逐一展开，并用 Google Brain 六年机器人研究史解释为什么团队最终选择 multi-task imitation。关键转变不是放弃在线学习，而是先让数据成为独立基础设施，使多个算法能够反复消费同一批真实经验。

\subsection{Ingredient 1：Scaling-Friendly Architecture}

第一类原料来自语言和视觉 scaling：高容量 network、self-attention、更多 compute/data、tokenization 和模块互操作。Transformer 的价值不只在 attention 公式，而是它提供一种统一序列接口，使图像历史、语言指令和动作都能进入同一 autoregressive model。

\lecturefigure{slide-07-ingredient-ml-scaling.jpg}{Ingredient 1：从 ML scaling 借鉴 architecture、compute、data 与 tokenization}{官方录像恢复幻灯片 V007；Stanford CS25 Lecture 15}

\begin{knowledgebox}{Tokenization 的机器人含义}
在 NLP 中 token 通常是离散子词；在 RT-1 中，图像先变成紧凑视觉 token，连续控制量被量化成 action tokens。统一 token interface 允许复用 Transformer，但量化误差、序列长度和实时解码成本会成为新的工程约束。
\end{knowledgebox}

\subsection{Ingredient 2：借用 Internet-Scale Models}

第二类原料不要求机器人团队从零训练语言或视觉语义。外部模型会在自身领域持续变好，可以提供 common sense、语言理解、视觉类别和可组合表示。对机器人而言，这些模型最适合做 planner、captioner、labeler 或 representation provider，而不是在没有 grounding 时直接输出 motor command。

\lecturefigure{slide-08-ingredient-internet-models.jpg}{Ingredient 2：利用不断进步的 Internet-scale generative models}{官方录像恢复幻灯片 V008；Stanford CS25 Lecture 15}

\teachervoice{讲者把这一点称为“Bitter Lesson 2.0”：不只选择会随 compute 变好的算法，还要把机器人接口设计成未来更强 foundation model 可以直接插入。今天的语言模型不是终点，系统应当吃到明天模型进步的红利。}

\subsection{Ingredient 3：Robotics 从 Online 转向 Offline}

第三类原料是数据制度。在线 RL 在机器人上需要边执行边学习，失败会损坏设备或浪费操作时间；offline learning 则先固定数据集，再用同一批 trajectories 训练和比较多个模型。slide 9 把 The Pile、LAION 等离线语料与机器人传统 online / on-policy 学习并置，强调数据消费方式的差异。

\lecturefigure{slide-09-ingredient-offline-data.jpg}{Ingredient 3：foundation models 消费巨型离线语料，而机器人长期依赖在线采集}{官方录像恢复幻灯片 V009；Stanford CS25 Lecture 15}

设离线示范数据为 $\mathcal{D}=\{(o_t,l,a_t)\}$，其中 $o_t$ 是观测、$l$ 是语言任务、$a_t$ 是专家动作。最直接的 behavioral cloning 目标是
\begin{equation}
\theta^*=\arg\min_{\theta}
-\mathbb{E}_{(o_t,l,a_t)\sim\mathcal{D}}
\log p_{\theta}(a_t\mid o_{\le t},l).
\end{equation}
它把控制学习化成监督学习，稳定而容易扩展；但 policy 一旦偏离专家轨迹，就会遇到训练分布没有覆盖的状态，这也是 offline imitation 的根本限制。

\subsection{六年历史：从 End-to-End 到 Multi-Task}

历史 slide 把 2016--2020 的 goal conditioning、QT-Opt、sim2real 和大规模真实机器人 RL，与 2020--2022 的 BC-Z、IL+RL、MT-Opt 等多任务系统连接起来。它不是完整综述，而是团队内部问题的变化：先问“端到端机器人能否学会一个任务”，再问“如何在更复杂场景中扩展到许多任务”。

\lecturefigure{slide-10-google-robotics-history.jpg}{Google Brain robotics 历史：从单任务端到端学习到多任务数据与策略}{官方录像恢复幻灯片 V010；Stanford CS25 Lecture 15}

\lecturefigure{slide-11-online-to-offline-shift.jpg}{六年转向：从 online methods 到 offline multi-task imitation}{官方录像恢复幻灯片 V011；Stanford CS25 Lecture 15}

\teachervoice{在学生询问“24/7 robot farm 是否足够”时，讲者强调机器人小时数不是唯一尺度。如果全天采集的都是同一类动作和场景，数据量增加却不扩展 task diversity；真正稀缺的是覆盖新的 failure mode 与组合变化。}

\begin{knowledgebox}{BC-Z 在课程中的位置}
BC-Z 使用语言或人类视频作为 task conditioning，在约百种 manipulation tasks 上训练 multi-task imitation policy，并测试 zero-shot task combinations。它证明 templated language 足以充当任务接口，也暴露 ResNet policy 对训练分布和新场景的敏感性，为 RT-1 提供数据与架构起点。
\end{knowledgebox}

\subsection{把 Ingredient 变成 Recipe}

最终配方把三列原料和 lessons 对齐：高容量 architecture / tokenization；可替换、持续变强的 external models；大量且多样的离线机器人数据。底部句子将它们压缩为“large diverse offline datasets + high-capacity architectures + language as universal glue”。

\lecturefigure{slide-12-foundation-recipe.jpg}{Turning Ingredients into a Recipe：架构、外部模型与离线数据的组合}{官方录像恢复幻灯片 V012；Stanford CS25 Lecture 15}

\begin{importantbox}{Data Generation 与 Data Consumption 解耦}
在线系统常为某个算法定制采集流程，算法改变就要重新跑机器人；offline recipe 把轨迹变成长期资产，使 RT-1、未来 RL 或不同 architecture 都能消费同一数据。解耦不会消除 distribution shift，却显著降低迭代成本并提高可复现性。
\end{importantbox}

\subsection{本章小结}

课堂 recipe 的核心不是“Transformer 万能”，而是统一接口与数据制度：tokenization 让多模态进入同一模型，外部 foundation models 提供可替换 prior，offline datasets 让采集和训练解耦。RT-1 接下来会把这些原则放入严格的实时控制预算。

